%% NOTATIONS -- original pages xvii-xviii
\fmhead[17]{NOTATIONS}{Notations}{notations}{NOTATIONS}

\runhead{General notations}

Positive means $\geq 0$.

$\Z$ (resp. $\Q$, $\R$, $\C$) is the ring (resp. the field) of integers
(resp. of rational numbers, of real numbers, of complex numbers).

If $p$ is a prime number, $\F_p$ denotes the prime field $\Z/p\Z$ and
$\Z_p$ (resp. $\Q_p$) the ring of $p$-adic integers (resp. the field of
$p$-adic rational numbers). One has:
\[
\Z_p=\varprojlim\nolimits_n \Z/p^n\Z\ ,\qquad \Q_p=\Z_p\big[\tfrac1p\big]\,.
\]

\runhead{Prime numbers}

They are denoted by $\ell$, $\ell'$, $p$, \dots; we mostly use the letter
$\ell$ for ``\textit{l}-adic representations'' and the letter $p$ for the
residue characteristic of some valuation.

\runhead{Fields}

If $K$ is a field, we denote by $\bar K$ an algebraic closure of $K$, and
by $K_s$ the separable closure of $K$ in $\bar K$; most of the fields we
consider are perfect, in which case $K_s=\bar K$.

If $L/K$ is a (possibly infinite) Galois extension, we denote its Galois
group by $\Gal(L/K)$; it is a projective limit of finite groups.

\runhead{Algebraic groups}

If $G$ is an algebraic group over a field $K$, and if $K'$ is a
commutative $K$-algebra, we denote by $G(K')$ the group of $K'$-points of
$G$ (the ``$K'$-rational'' points of $G$).

When $K'$ is a field, we denote by $G/K'$ the $K'$-algebraic group
$G\times_K K'$ obtained from $G$ by extending the ground field from $K$
to $K'$.

Let $V$ be a finite dimensional $K$-vector space. We denote by
$\Aut_K(V)$, or $\Aut(V)$, the group of its $K$-linear automorphisms, and
by $\GLV$ the corresponding $K$-algebraic group (cf.~chap.~\ref{chap:I},
\S\ref{sec:I.2.4}).

For any commutative $K$-algebra $K'$, the group $\GLV(K')$ of
$K'$-points of $\GLV$ is $\Aut_{K'}(V\tens_K K')$; for instance,
$\GLV(K)=\Aut(V)$.
